%% ch04.tex --- Multiple Linear Regression and Vectorization
%% Numbers are computed by verify/ch04_multiple_features.py.

\chapter{Multiple Linear Regression and Vectorization}
\label{ch:ch04}

\epigraph{Civilisation advances by extending the number of important operations
which we can perform without thinking about them.}%
{---\ Alfred North Whitehead}

\begin{chapterlearning}
After working through this chapter you will be able to:
\begin{enumerate}
  \item Use the notation of multiple features correctly, including the design
        matrix and the distinction between $\xii{i}$ and $\xjj{i}{j}$.
  \item Write the model, the cost and the gradient in vector and matrix form,
        and say what shape every quantity has.
  \item Explain what vectorization does and does not change, and estimate the
        speed-up it buys.
  \item Derive $\partial J/\partial w_j$ for $n$ features and implement gradient
        descent without a loop over examples or features.
  \item Derive the normal equation from the orthogonality of the residual, and
        state its computational cost.
  \item Choose between gradient descent and the normal equation for a given
        problem, with reasons.
  \item Recognise when $\Xmat^{\mathsf{T}}\Xmat$ is singular or nearly so, say
        what each situation means, and know what a solver returns anyway.
  \item Explain why a fitted coefficient in a multi-feature model is not the
        effect of that feature.
\end{enumerate}
\end{chapterlearning}

\section{More Than One Feature}
\label{sec:more-features}

Real prediction problems rarely offer a single input. The apartment price of
\cref{ch:ch01} depends on floor area, but also on the age of the building, the
number of rooms, the floor it is on, and the distance to the nearest station.
Nothing about \cref{ch:ch02,ch:ch03} required the input to be one number; it
merely made the pictures drawable. We now drop that restriction.

Write $n$ for the number of features. The $i$-th training example is now a
vector, and we need a second index to reach inside it.

\begin{defbox}{Notation for multiple features}
\begin{center}
\small
\renewcommand{\arraystretch}{1.3}
\begin{tabular}{@{}lL{8.2cm}@{}}
\toprule
$m$ & number of training examples\\
$n$ & number of features\\
$\vx^{(i)}$ & the $i$-th example, a vector in $\Real{n}$\\
$\xjj{i}{j}$ & feature $j$ of example $i$ --- a single number\\
$\Xmat$ & the \term{design matrix}: $m\times n$, one row per example, one
          column per feature\\
$\vy$ & the vector of targets, in $\Real{m}$\\
$\vw$ & the weight vector, in $\Real{n}$; $b$ remains a single number\\
\bottomrule
\end{tabular}
\end{center}
Rows are examples and columns are features. This convention is universal in
machine learning libraries and is worth committing to memory, because half of
all shape errors come from silently transposing it.
\end{defbox}

The running example for this chapter is five apartments described by two
features: floor area in hundreds of square feet, and the age of the building in
years. It is dataset~D2 of \cref{ch:appD}.

\begin{center}
\small
\renewcommand{\arraystretch}{1.2}
\begin{tabular}{@{}crrr@{}}
\toprule
$i$ & area $\xjj{i}{1}$ & age $\xjj{i}{2}$ & price $\yii{i}$\\
\midrule
$1$ & $8$  & $20$ & $30$\\
$2$ & $12$ & $10$ & $46$\\
$3$ & $15$ & $12$ & $55$\\
$4$ & $18$ & $5$  & $72$\\
$5$ & $22$ & $8$  & $84$\\
\bottomrule
\end{tabular}
\end{center}

\noindent So $m=5$, $n=2$, $\vx^{(3)} = (15, 12)$ and $\xjj{3}{2} = 12$.

\begin{defbox}{Multiple linear regression}
With $n$ features the model is
\[
  f_{\vw,b}(\vx) \;=\; w_1x_1 + w_2x_2 + \dots + w_nx_n + b
                \;=\; \vw\cdot\vx + b ,
\]
with $n+1$ parameters: the weight vector $\vw\in\Real{n}$ and the bias $b$.
Geometrically it is a hyperplane in $n+1$ dimensions --- a line when $n=1$, a
plane when $n=2$.
\end{defbox}

\begin{figure}[htbp]
  \centering
  \begin{tikzpicture}
    \begin{axis}[width=0.78\linewidth, height=6.0cm,
      view={125}{22},
      xlabel={area}, ylabel={age}, zlabel={price},
      tick label style={font=\tiny}, label style={font=\small},
      colormap={ml}{color=(MLLight) color=(MLTeal)},
      samples=2, domain=7:23, y domain=3:21,
      zmin=15, zmax=95]
      \addplot3[surf,shader=interp,opacity=0.55,forget plot]
        {7.093085 + 3.620567*x - 0.363780*y};
      \addplot3[only marks,mark=*,mark size=2pt,MLRose] coordinates {
        (8,20,30)(12,10,46)(15,12,55)(18,5,72)(22,8,84)};
    \end{axis}
  \end{tikzpicture}
  \caption{Two features, one target. The model is a plane; the five training
  examples sit near but not on it. The cost is still the mean squared vertical
  gap, and gradient descent still walks downhill --- in a three-dimensional
  parameter space now, which we can no longer draw.}
  \label{fig:plane}
\end{figure}

\section{The Model in Vector Form}
\label{sec:vector-form}

The sum $w_1x_1 + \dots + w_nx_n$ is a dot product, and writing it as one is not
mere shorthand: it is what makes the next section's speed-up possible and the
one after that's derivation short.

For all $m$ predictions at once, stack the examples as the rows of $\Xmat$:
\[
  \hat{\vy} \;=\; \Xmat\vw + b\mathbf{1},
  \qquad
  \underbrace{\Xmat}_{m\times n}\ \underbrace{\vw}_{n\times 1}
  \;\to\;\underbrace{\hat{\vy}}_{m\times 1}.
\]
Checking shapes this way, every time, prevents most of the errors you will
otherwise make.

\begin{figure}[htbp]
  \centering
  \begin{tikzpicture}[every node/.style={font=\scriptsize}]
    \matrix (X) [matrix of math nodes, left delimiter=(, right delimiter=),
                 row sep=1pt, column sep=6pt, nodes={minimum width=1.1cm}]
    {
      8  & 20\\
      12 & 10\\
      15 & 12\\
      18 & 5\\
      22 & 8\\
    };
    \node[above=1pt of X, font=\scriptsize\bfseries] {$\Xmat$ \ ($5\times 2$)};
    \draw[decorate,decoration={brace,amplitude=5pt},MLTeal]
      ($(X.south west)+(-0.30,0)$) -- ($(X.north west)+(-0.30,0)$)
      node[midway,left=9pt,MLTeal,align=right] {$m=5$\\ examples};
    \draw[decorate,decoration={brace,amplitude=4pt,mirror},MLAmber]
      ($(X.south west)+(0,-0.25)$) -- ($(X.south east)+(0,-0.25)$)
      node[midway,below=4pt,MLAmber] {$n=2$ features};
    \node[right=1.4cm of X] (w)
      {$\begin{pmatrix} w_1\\ w_2\end{pmatrix}$};
    \node[right=0.3cm of w] (plus) {$+\;b\mathbf{1}\;=$};
    \node[right=0.2cm of plus] (yh)
      {$\begin{pmatrix}\hat y^{(1)}\\ \hat y^{(2)}\\ \hat y^{(3)}\\
                       \hat y^{(4)}\\ \hat y^{(5)}\end{pmatrix}$};
    \node[above=1pt of w, font=\scriptsize\bfseries] {$\vw$};
    \node[above=1pt of yh, font=\scriptsize\bfseries] {$\hat{\vy}$};
  \end{tikzpicture}
  \caption{One matrix--vector product computes every prediction. Rows are
  examples; columns are features. The bias is added to all $m$ entries by
  broadcasting.}
  \label{fig:design-matrix}
\end{figure}

It is often convenient to absorb $b$ into $\vw$ by prepending a column of ones
to $\Xmat$ and a component $w_0 = b$ to the weight vector. Then
$\hat{\vy} = \Xmat\bm{\theta}$ with $\bm{\theta}\in\Real{n+1}$, and the bias
needs no special treatment. We use this in \cref{sec:normal-equation}; the code
in between keeps $b$ separate, because that is what most libraries do.

\section{Vectorization}
\label{sec:vectorization}

Consider three ways to compute one prediction.

\begin{lstlisting}[style=pythonml,caption={The same dot product, three ways.
Only the third is acceptable at scale.}]
# 1. explicit loop -- correct, and unusably slow
f = 0.0
for j in range(n):
    f += w[j] * x[j]
f += b

# 2. the same loop, written as a comprehension -- no faster
f = sum(w[j] * x[j] for j in range(n)) + b

# 3. vectorized -- one call into a compiled kernel
f = np.dot(w, x) + b          # or:  w @ x + b
\end{lstlisting}

All three compute the same number. The third is faster by a factor that is worth
measuring rather than guessing:

\begin{center}
\small
\renewcommand{\arraystretch}{1.2}
\begin{tabular}{@{}rrrr@{}}
\toprule
vector length $N$ & Python loop & \code{numpy} & speed-up\\
\midrule
$1{,}000$     & $0.13$ ms   & $0.0010$ ms & $136\times$\\
$100{,}000$   & $13.29$ ms  & $0.0029$ ms & $4521\times$\\
$1{,}000{,}000$ & $139.88$ ms & $0.1661$ ms & $842\times$\\
\bottomrule
\end{tabular}
\end{center}

\noindent (Measured on the machine used to prepare this book; your factors will
differ, but not by so much that the conclusion changes. The results agree to
within $2\times 10^{-9}$, the accumulated difference between summing in a
different order.)

Notice that the speed-up is not monotone in $N$. At $N = 100{,}000$ both vectors
fit in fast cache and the vectorized version is at its most impressive; by
$N = 10^6$ it is limited by how fast memory can be read, and the advantage
falls back. Performance work is full of such non-monotonicities, which is
precisely why the numbers above were measured and not estimated.

Four things account for the difference, and none of them is a better algorithm:

\begin{enumerate}
  \item \textbf{No interpreter overhead.} The Python loop executes several
        bytecode instructions and allocates an object per element. The
        vectorized call enters compiled code once.
  \item \textbf{SIMD instructions.} A modern core multiplies four or eight
        double-precision numbers per instruction. A loop written element by
        element cannot use them.
  \item \textbf{Memory locality.} Contiguous arrays are prefetched; Python lists
        of boxed floats are scattered across the heap.
  \item \textbf{Parallelism.} Larger matrix operations are handed to a
        multi-threaded BLAS library, and on a GPU to thousands of cores at once.
\end{enumerate}

\begin{keyinsight}{Vectorization is not an optimization of the algorithm}
The vectorized code performs the same arithmetic, in the same quantity, to the
same accuracy. What changes is who executes it. This is why vectorizing is
almost always worth doing and almost never changes an answer --- and why, if
vectorizing your code \emph{does} change an answer by more than rounding, you
have a bug, not a speed-up.
\end{keyinsight}

\begin{pitfall}{Shapes: $(m,)$ is not $(m,1)$}
The commonest failure in a first vectorized implementation is not slowness but
silent broadcasting. In NumPy, an array of shape $(m,)$ and one of shape
$(m,1)$ are different objects, and subtracting the second from the first
produces an $(m,m)$ matrix rather than an error:

\begin{lstlisting}[style=pythonml]
y     = np.array([30., 46., 55., 72., 84.])  # (5,)
y_hat = (X @ w + b).reshape(-1, 1)           # (5,1) oops
r     = y_hat - y                            # (5,5)!
J     = (r ** 2).sum() / (2 * m)             # wrong
\end{lstlisting}

Line 2 introduces the mistake; line 3 is where it does its damage, subtracting
a $(5,)$ array from a $(5,1)$ one to produce a $(5,5)$ matrix of every
difference against every other. The cost on line 4 comes out finite, positive
and meaningless, and gradient descent runs happily on it. Print \code{.shape} at every step while developing, and assert
the shapes you expect:
\code{assert r.shape == (m,)}.
\end{pitfall}

\section{Gradient Descent for \texorpdfstring{$n$}{n} Features}
\label{sec:gd-multi}

The cost is unchanged in form:
\[
  J(\vw, b) \;=\; \frac{1}{2m}\sum_{i=1}^{m}
    \bigl(f_{\vw,b}(\vx^{(i)}) - \yii{i}\bigr)^{2} .
\]

\begin{derivation}{The gradient with respect to each weight}
Differentiate with respect to one weight $w_j$. Only the $j$-th term of the dot
product depends on it, so the inner derivative is $\xjj{i}{j}$:
\[
  \frac{\partial J}{\partial w_j}
    = \frac{1}{m}\sum_{i=1}^{m}
      \bigl(f_{\vw,b}(\vx^{(i)}) - \yii{i}\bigr)\,\xjj{i}{j},
  \qquad j = 1,\dots,n,
\]
and, as before, $\partial J/\partial b$ is the same expression with the weight
$\xjj{i}{j}$ replaced by $1$. Compare \cref{sec:gd-linreg}: the univariate
formula is the case $n=1$, and nothing else has changed.

Collecting the $n$ partial derivatives into a vector and writing
$\mathbf{r} = \Xmat\vw + b\mathbf{1} - \vy$ for the residual vector,
\begin{equation}
  \grad_{\vw} J \;=\; \frac{1}{m}\,\Xmat^{\mathsf{T}}\mathbf{r},
  \qquad
  \frac{\partial J}{\partial b} \;=\; \frac{1}{m}\,\mathbf{1}^{\mathsf{T}}\mathbf{r},
  \label{eq:vector-gradient}
\end{equation}
because the $j$-th component of $\Xmat^{\mathsf{T}}\mathbf{r}$ is exactly
$\sum_i \xjj{i}{j} r^{(i)}$ --- the $j$-th column of $\Xmat$ dotted with the
residuals.
\end{derivation}

\Cref{eq:vector-gradient} is one of the most useful formulas in this book. Two
matrix products compute all $n+1$ derivatives: $\Xmat\vw$ forward to get the
residuals, $\Xmat^{\mathsf{T}}\mathbf{r}$ backward to get the gradient. That
forward-then-backward structure is not a coincidence, and \cref{ch:ch09} names
it backpropagation.

\begin{algorithmbox}{Batch gradient descent for multiple linear regression}
\textbf{Input:} $\Xmat$ ($m\times n$), $\vy$ ($m$), learning rate $\alpha$,
iterations $N$.

\begin{enumerate}\setlength{\itemsep}{1pt}
  \item $\vw \leftarrow \mathbf{0}$, \quad $b \leftarrow 0$.
  \item \textbf{for} $k = 1,\dots,N$:
  \begin{enumerate}\setlength{\itemsep}{1pt}
    \item $\mathbf{r} \leftarrow \Xmat\vw + b\mathbf{1} - \vy$
          \hfill \emph{($O(mn)$)}
    \item $\mathbf{g} \leftarrow \frac{1}{m}\Xmat^{\mathsf{T}}\mathbf{r}$,
          \quad $g_b \leftarrow \frac{1}{m}\sum_i r^{(i)}$
          \hfill \emph{($O(mn)$)}
    \item $\vw \leftarrow \vw - \alpha\mathbf{g}$,
          \quad $b \leftarrow b - \alpha g_b$ \hfill \emph{(simultaneously)}
  \end{enumerate}
\end{enumerate}

\textbf{Cost:} $O(mn)$ per iteration, $O(mnN)$ in total --- linear in the number
of features.
\end{algorithmbox}

\begin{lstlisting}[style=pythonml,caption={The whole algorithm. There is no loop
over examples and no loop over features.}]
def gradient_descent(X, y, alpha, iters):
    m, n = X.shape
    w, b = np.zeros(n), 0.0
    for _ in range(iters):
        r  = X @ w + b - y                 # (m,)  residuals
        gw = (X.T @ r) / m                 # (n,)  all n gradients at once
        gb = r.mean()                      # scalar
        w, b = w - alpha * gw, b - alpha * gb    # simultaneous
    return w, b
\end{lstlisting}

\section{The Normal Equation}
\label{sec:normal-equation}

For this model the exact solution is again available, and the derivation is
shorter in matrix form than the one-variable version of \cref{sec:closed-form}
--- and more illuminating, because it is the projection argument of
\cref{sec:projection} with $n+1$ spanning vectors instead of two.

\begin{derivation}{The normal equation}
Absorb the bias: let $\Xmat$ now include a leading column of ones, so
$\bm{\theta}\in\Real{n+1}$ and $\hat{\vy} = \Xmat\bm{\theta}$. As
$\bm{\theta}$ ranges over $\Real{n+1}$, $\hat{\vy}$ sweeps the column space of
$\Xmat$ --- the subspace of $\Real{m}$ spanned by the columns.

Minimising $J = \frac{1}{2m}\|\Xmat\bm{\theta}-\vy\|^2$ means finding the point
of that subspace closest to $\vy$, which is its orthogonal projection. The
residual $\mathbf{r} = \Xmat\bm{\theta}-\vy$ must therefore be perpendicular to
\emph{every} column of $\Xmat$, which is $n+1$ equations written at once as
\begin{equation}
  \Xmat^{\mathsf{T}}\bigl(\Xmat\bm{\theta} - \vy\bigr) = \mathbf{0}
  \qquad\Longleftrightarrow\qquad
  \Xmat^{\mathsf{T}}\Xmat\,\bm{\theta} = \Xmat^{\mathsf{T}}\vy ,
  \label{eq:normal-equations}
\end{equation}
the \term{normal equations} --- ``normal'' in the sense of perpendicular. When
$\Xmat^{\mathsf{T}}\Xmat$ is invertible,
\begin{equation}
  \bm{\theta}^\star
    = \bigl(\Xmat^{\mathsf{T}}\Xmat\bigr)^{-1}\Xmat^{\mathsf{T}}\vy .
  \label{eq:normal-equation}
\end{equation}
Setting $\grad J = \frac1m \Xmat^{\mathsf{T}}\mathbf{r} = \mathbf{0}$ in
\cref{eq:vector-gradient} gives the same equations, as it must: the stationarity
condition and the perpendicularity condition are one statement.
\end{derivation}

\begin{workedexample}{Dataset D2, both ways}
With the bias column, $\Xmat$ is $5\times 3$ and
\[
  \Xmat^{\mathsf{T}}\Xmat =
  \begin{pmatrix}
    5    & 75   & 55\\
    75   & 1241 & 726\\
    55   & 726  & 733
  \end{pmatrix},
  \qquad
  \Xmat^{\mathsf{T}}\vy =
  \begin{pmatrix} 287\\ 4761\\ 2752 \end{pmatrix}.
\]
The top-left entry is $m$; the rest of the first row and column are the feature
sums; the interior is the matrix of feature cross-products. Solving the
$3\times3$ system:
\[
  b = 7.093085,
  \qquad
  w_1 = 3.620567,
  \qquad
  w_2 = -0.363780 ,
\]
with $J = 0.889051$. The signs are as one would hope: price rises with area and
falls with age.

\textbf{The same answer by gradient descent.} Running the algorithm of
\cref{sec:gd-multi} on standardised features (\cref{ch:ch05} explains why that
is necessary here), for increasing budgets:

\begin{center}
\small
\renewcommand{\arraystretch}{1.1}
\begin{tabular}{@{}rrrrr@{}}
\toprule
iterations & $b$ & $w_1$ & $w_2$ & $J$\\
\midrule
$10$      & $28.589218$ & $2.712345$ & $-1.226892$ & $5.786337$\\
$100$     & $7.219119$  & $3.615619$ & $-0.368491$ & $0.889157$\\
$1{,}000$ & $7.093085$  & $3.620567$ & $-0.363780$ & $0.889051$\\
\midrule
normal eq.& $7.093085$  & $3.620567$ & $-0.363780$ & $0.889051$\\
\bottomrule
\end{tabular}
\end{center}

After a thousand iterations the two agree to within $1.4\times 10^{-13}$ --- the
limit of double precision, not of the algorithm. Two entirely different methods,
one answer: this is the check to run whenever you implement either.

\textbf{Is the second feature worth having?} Fitting area alone gives
$b = -1.5655$, $w_1 = 3.9310$ and $J = 1.4648$. Adding age reduces the cost to
$0.8891$, a $39\%$ reduction. Note that $w_1$ also changed, from $3.93$ to
$3.62$: the coefficient on area is not a property of area alone.
\end{workedexample}

\subsection{What it costs}

Forming $\Xmat^{\mathsf{T}}\Xmat$ takes $O(mn^2)$ operations and solving the
resulting system takes $O(n^3)$. The cubic term is what eventually stops you:

\begin{center}
\small
\renewcommand{\arraystretch}{1.15}
\begin{tabular}{@{}rrr@{}}
\toprule
$n$ & time to solve & ratio to $n/2$\\
\midrule
$200$   & $0.108$ ms   & ---\\
$400$   & $0.491$ ms   & $4.5\times$\\
$800$   & $2.306$ ms   & $4.7\times$\\
$1{,}600$ & $16.440$ ms & $7.1\times$\\
$3{,}200$ & $129.086$ ms & $7.9\times$\\
\bottomrule
\end{tabular}
\end{center}

Each doubling of $n$ should multiply the time by $2^3 = 8$, and the measured
ratios approach that once the matrices outgrow the processor's cache. Extrapolate:
at $n = 100{,}000$ features --- an ordinary size for text --- the same solve would
take roughly $10^{4.5}$ times as long as at $n=3{,}200$, which is several days.
Gradient descent, at $O(mn)$ per iteration, is unbothered.

\begin{center}
\small
\renewcommand{\arraystretch}{1.25}
\begin{tabular}{@{}L{5.1cm}L{5.1cm}@{}}
\toprule
\textbf{Normal equation} & \textbf{Gradient descent}\\
\midrule
No learning rate to choose. & $\alpha$ must be chosen, and badly chosen it
  fails.\\
No iterations, no stopping rule. & Needs a budget and a stopping rule.\\
No feature scaling required. & Feature scaling matters a great deal
  (\cref{ch:ch05}).\\
$O(n^3)$: impractical beyond roughly $10^4$ features. & $O(mn)$ per iteration:
  fine at any $n$.\\
Fails outright if $\Xmat^{\mathsf{T}}\Xmat$ is singular. & Unbothered by
  singularity; returns \emph{a} minimiser.\\
Exists only for linear regression with squared error. & Works for every model in
  this book.\\
\bottomrule
\end{tabular}
\end{center}

The practical rule: with a few thousand features or fewer and a squared-error
linear model, use the normal equation --- it is exact and needs no tuning.
Otherwise, or as soon as the model stops being linear, use gradient descent.

\section{When \texorpdfstring{$\Xmat^{\mathsf{T}}\Xmat$}{X'X} Is Not Invertible}
\label{sec:singular}

\Cref{eq:normal-equation} assumes an inverse exists. Two situations in which it
does not are common enough to recognise on sight.

\textbf{Redundant features.} Suppose someone records floor area twice, once in
square feet and once in square metres, or --- as happens more often --- a feature
is accidentally duplicated by a join. Then one column of $\Xmat$ is a multiple of
another, the columns are linearly dependent, and $\Xmat^{\mathsf{T}}\Xmat$ is
singular. Duplicating the area column of D2 gives a $4\times 4$ matrix
$\Xmat^{\mathsf{T}}\Xmat$ with determinant exactly $0$ and rank $3$.

\textbf{More features than examples.} If $n \ge m$ then $\Xmat$ has at most $m$
independent columns, so $\Xmat^{\mathsf{T}}\Xmat$ ($n\times n$) cannot have full
rank. There are then infinitely many parameter vectors achieving $J = 0$, all of
them fitting the training data perfectly and most of them useless ---
\cref{ch:ch07}'s overfitting in its purest form.

\begin{pitfall}{Your solver will not tell you}
Calling \code{np.linalg.inv} on a singular matrix raises an error, but almost
nobody calls it: the usual routines, \code{np.linalg.lstsq} and
\code{sklearn.linear\_model.LinearRegression}, use the pseudo-inverse and return
an answer without complaint. For the duplicated-area version of D2 the solver
returns
\[
  b = 7.0931,\quad w_1 = 1.8103,\quad w_2 = -0.3638,\quad w_3 = 1.8103,
\]
where $w_1$ and $w_3$ are the coefficients on the two identical columns. It has
split the original coefficient exactly in half: $1.8103 + 1.8103 = 3.620567$,
the value from the non-degenerate fit. This is the \emph{minimum-norm} solution
--- of the infinitely many equally good answers, the one with the smallest
$\|\bm{\theta}\|$.

The predictions are perfect and identical to before. The coefficients are
arbitrary. Anyone reading them as ``the effect of area'' would be reading half
of it, twice, and would draw the wrong conclusion. Check for duplicated and
derived columns before fitting, and check the rank of $\Xmat$ when a coefficient
surprises you.
\end{pitfall}

The remedies, in order of preference: remove the redundant feature; collect more
examples if $n \ge m$; or add regularization, which makes
$\Xmat^{\mathsf{T}}\Xmat + \lambda I$ invertible for any $\lambda>0$ and is the
subject of \cref{ch:ch07}.

\section{Multicollinearity, and Reading Coefficients}
\label{sec:multicollinearity}

Exact dependence between features is easy to spot. Near-dependence is not, and it
is far more common: floor area and number of rooms, income and years of
education, height and weight all move together without being multiples of one
another.

Replace D2's age column with a feature that is almost exactly three times the
area, differing only in the second decimal place. The matrix
$\Xmat^{\mathsf{T}}\Xmat$ is now invertible, so everything appears to work, and
its \term{condition number} --- the ratio of its largest eigenvalue to its
smallest --- is $1.3\times 10^{8}$. The fitted coefficients are

\[
  b = 0.401, \qquad w_1 = 306.963, \qquad w_2 = -101.050 .
\]

These are absurd. The sensible single-feature answer was $w_1 = 3.93$; the model
has found a pair of enormous coefficients that nearly cancel, since
$306.963 - 3\times 101.050 \approx 3.81$. The predictions are perfectly
reasonable; only the individual coefficients are nonsense.

Worse, they are unstable. Add $0.1$ to a single target value --- a rounding
difference in one price --- and refit:
\[
  b = 0.491, \qquad w_1 = 303.920, \qquad w_2 = -100.037 .
\]
$w_1$ has moved by $3.04$, which is nearly the size of the entire meaningful
coefficient from the single-feature fit. A dataset that differs by one part in
a thousand yields a materially different model.

\begin{pitfall}{A coefficient is not the effect of a feature}
The standard gloss --- ``$w_j$ is the change in $y$ per unit change in $x_j$,
holding the other features fixed'' --- is mathematically correct and practically
treacherous, because it presumes the other features \emph{can} be held fixed. If
$x_2$ is always about $3x_1$, then no apartment exists with $x_1$ increased and
$x_2$ unchanged, and the coefficient describes a comparison the data never made.

Symptoms of multicollinearity: coefficients far larger than the target's scale;
signs that contradict domain knowledge; coefficients that swing wildly when a
single example is added or removed. Diagnose it with the condition number of
$\Xmat^{\mathsf{T}}\Xmat$ or a correlation matrix of the features. Treat it by
removing one of the pair, combining them into a single feature, or regularizing
(\cref{ch:ch07}). And be clear about your goal: if you only need predictions,
multicollinearity is nearly harmless; if you intend to \emph{interpret} the
coefficients, it is fatal.
\end{pitfall}

One more warning about reading coefficients, which does not require any
degeneracy at all. In the honest D2 fit, $w_1 = 3.62$ and $w_2 = -0.36$. It is
tempting to conclude that area matters ten times as much as age. It does not
follow: $w_1$ is measured per hundred square feet and $w_2$ per year, and the two
are not comparable any more than kilograms and kilometres are. Comparing the
importance of features requires putting them on a common scale first, which is
the first thing \cref{ch:ch05} does.

\begin{labbox}{Lab 4 --- Vectorized multiple regression}
Load dataset D2 and implement \code{predict}, \code{cost} and
\code{gradient\_descent} with no loop over examples or features. Confirm
$J = 0.889051$ at the fitted parameters. Implement the normal equation with
\code{np.linalg.solve} --- not \code{inv} --- and check that the two methods
agree to at least ten decimal places. Time your loop-based and vectorized cost
functions on a random problem with $m = 10^5$ and $n = 100$, and report the
ratio. Finally, duplicate the area column, refit, and confirm that the
coefficient splits in half while the predictions do not move.
\end{labbox}

\section*{Chapter Summary}
\addcontentsline{toc}{section}{Chapter Summary}

\begin{itemize}
  \item With $n$ features the model is $f_{\vw,b}(\vx) = \vw\cdot\vx + b$. The
        design matrix $\Xmat$ is $m\times n$: rows are examples, columns are
        features.
  \item All $m$ predictions are one matrix--vector product,
        $\hat{\vy} = \Xmat\vw + b\mathbf{1}$. Checking shapes at every step
        prevents most implementation errors; $(m,)$ and $(m,1)$ are different,
        and mixing them broadcasts silently instead of failing.
  \item Vectorization performs identical arithmetic in compiled, SIMD,
        cache-friendly, sometimes parallel code. Measured speed-ups here ranged
        from $136\times$ to $4521\times$. It changes the implementation, never
        the algorithm or the answer.
  \item $\partial J/\partial w_j = \frac1m\sum_i r^{(i)}\xjj{i}{j}$, or in
        vector form $\grad_{\vw}J = \frac1m\Xmat^{\mathsf{T}}\mathbf{r}$: one
        product forward for the residuals, one backward for the gradient. The
        update rule of \cref{ch:ch03} is otherwise unchanged.
  \item The normal equations $\Xmat^{\mathsf{T}}\Xmat\bm{\theta} =
        \Xmat^{\mathsf{T}}\vy$ state that the residual is perpendicular to every
        column of $\Xmat$ --- the projection picture of \cref{ch:ch02} in $n+1$
        dimensions. They cost $O(n^3)$ to solve.
  \item Use the normal equation for a linear model with up to a few thousand
        features; use gradient descent otherwise, and always once the model
        stops being linear.
  \item $\Xmat^{\mathsf{T}}\Xmat$ is singular when features are redundant or
        $n \ge m$. Solvers return the minimum-norm answer rather than an error,
        so the failure is silent and shows up only in the coefficients.
  \item Near-collinear features give huge, unstable, uninterpretable
        coefficients while leaving predictions intact. Coefficients of different
        units are not comparable in size.
\end{itemize}

\section*{Exercises}
\addcontentsline{toc}{section}{Exercises}

\exercisesection{Short Questions}

\sq For a dataset with $m = 400$ examples and $n = 12$ features, give the shape
of $\Xmat$, $\vw$, $\vy$ and $\Xmat^{\mathsf{T}}\Xmat$.

\sq What is the difference between $\vx^{(3)}$, $\xjj{3}{2}$ and $x_2$?

\sq Write the multiple linear regression model in both summation and vector
notation.

\sq Give two reasons a vectorized dot product is faster than the equivalent
Python loop.

\sq Does vectorizing a computation change its result? Explain your answer,
including any exception.

\sq Write $\grad_{\vw}J$ in matrix form and state the shape of each factor.

\sq State the normal equations and explain what the word ``normal'' refers to.

\sq Give the computational cost of the normal equation and of one iteration of
gradient descent, in terms of $m$ and $n$.

\sq Name two situations in which $\Xmat^{\mathsf{T}}\Xmat$ has no inverse, and
say what a typical solver does in each case.

\sq Why can two fitted coefficients not be compared by magnitude?

\exercisesection{Analytical Problems}

\ap For dataset D2, verify by hand that the $(2,3)$ entry of
$\Xmat^{\mathsf{T}}\Xmat$ is $726$ and that the second entry of
$\Xmat^{\mathsf{T}}\vy$ is $4761$. State in words what each of these sums
computes.

\ap Show that the $j$-th component of $\Xmat^{\mathsf{T}}\mathbf{r}$ equals
$\sum_{i=1}^{m} \xjj{i}{j}\,r^{(i)}$, and hence that
\cref{eq:vector-gradient} agrees with the component-wise derivative.

\ap Starting from $J = \frac{1}{2m}\|\Xmat\bm{\theta}-\vy\|^2$, expand the norm
and differentiate with respect to $\bm{\theta}$ to obtain
\cref{eq:normal-equations} directly, without the projection argument.

\ap Suppose feature $1$ is measured in metres and is then re-expressed in
centimetres. How do $w_1$, $b$, the fitted predictions and $J$ change? Now
suppose instead that the \emph{target} is re-expressed in cents rather than
dollars. Answer the same four questions.

\ap Show that if two columns of $\Xmat$ are identical, then for any solution
$\bm{\theta}$ there are infinitely many others with the same cost, obtained by
moving weight from one column to the other. Which of them does the pseudo-inverse
return, and why?

\ap A dataset has $m = 50$ examples and $n = 60$ features. Explain why a
zero-cost fit exists, why it is worthless, and what the smallest change to the
setup would be that makes the problem well posed.

\exercisesection{Design Problems}

\dpr You are given $80$ candidate features for a price model, many derived from
one another (area, area per room, log area, rooms, rooms per floor). Describe how
you would decide which to include, what diagnostics you would compute before
fitting, and how your answer would differ if the goal were prediction rather than
explanation.

\dpr A stakeholder asks: ``the model says each extra year of age reduces the
price by \$364 --- should we renovate?'' Write the reply, covering what the
coefficient does and does not license them to conclude, and what evidence would
actually answer their question.

\dpr Your training set has $m = 10^{7}$ examples and $n = 300$ features and does
not fit in memory. Choose between the normal equation and gradient descent,
justify the choice on cost grounds, and describe how you would compute what your
chosen method needs without ever holding $\Xmat$ in memory at once.

\dpr Design a set of automatic checks to run on $\Xmat$ before any model is
fitted --- rank, condition number, duplicated columns, constant columns, scale
--- and say for each what threshold would trigger a warning and what the warning
should tell the user to do.

\exercisesection{Programming Exercises}

\pe Implement \code{cost(X, y, w, b)} for the multi-feature model with no
explicit loop, and verify $J = 0.889051$ on D2 at the fitted parameters. Write a
loop-based version and \code{assert np.allclose} between them.

\pe Reproduce the timing table of \cref{sec:vectorization} on your own machine
for $N \in \{10^3, 10^5, 10^6\}$. Do you see the same non-monotone speed-up?
Explain what your numbers suggest about your machine's cache.

\pe Implement gradient descent from \cref{sec:gd-multi} and the normal equation
from \cref{sec:normal-equation}, and plot the maximum absolute difference between
their parameter vectors as a function of the number of iterations, on a
logarithmic axis. How many iterations are needed for agreement to six decimal
places?

\pe Generate a dataset with two nearly-collinear features and fit it $100$ times,
each time perturbing a single target by a small random amount. Plot the
distribution of $w_1$. Repeat with the collinearity removed and compare the two
distributions. This is the experiment behind \cref{sec:multicollinearity}.

\exercisesection{Self-Test}

\qz In the design matrix $\Xmat$, the rows correspond to: (a)~features;
(b)~examples; (c)~parameters; (d)~predictions.

\qz $\xjj{4}{2}$ denotes: (a)~example $2$ of feature $4$; (b)~feature $2$ of
example $4$; (c)~$x_4$ squared; (d)~the fourth parameter.

\qz Vectorizing a correct loop-based implementation should change the result by:
(a)~nothing at all; (b)~at most rounding error; (c)~a factor of $m$;
(d)~an unpredictable amount.

\qz The vector of all $n$ weight gradients equals: (a)~$\Xmat\mathbf{r}/m$;
(b)~$\Xmat^{\mathsf{T}}\mathbf{r}/m$; (c)~$\mathbf{r}^{\mathsf{T}}\Xmat/n$;
(d)~$\Xmat^{\mathsf{T}}\Xmat\bm{\theta}$.

\qz The normal equations express the fact that the residual is: (a)~zero;
(b)~minimal in absolute value; (c)~perpendicular to every column of $\Xmat$;
(d)~normally distributed.

\qz Solving the normal equations costs about: (a)~$O(mn)$; (b)~$O(n^2)$;
(c)~$O(n^3)$; (d)~$O(m^3)$.

\qz If two columns of $\Xmat$ are identical, \code{np.linalg.lstsq}:
(a)~raises an error; (b)~returns the minimum-norm solution; (c)~returns
infinity; (d)~drops one column automatically.

\qz Near-collinear features chiefly damage: (a)~the predictions;
(b)~the interpretability and stability of the coefficients; (c)~the cost
function; (d)~the learning rate.

\furtherreading

\textcite{james2021} §3.2 is the natural companion to this chapter: it develops
multiple linear regression, the interpretation of coefficients as partial
effects, and --- in §3.3.3 --- collinearity and the variance inflation factor,
which is the standard diagnostic for \cref{sec:multicollinearity}.
\textcite{bishop2006} §3.1 gives the same material in matrix form with the
geometry of \cref{sec:normal-equation} drawn explicitly, and introduces the
pseudo-inverse in the context of the singular case.

\bigskip
\noindent\emph{Next:} \Cref{ch:ch05} takes up the loose ends this chapter left
behind. Why did the worked example have to standardise the features before
gradient descent would work at all? How should $\alpha$ be chosen? And what do
you do when the relationship is not a straight line?
